\honest{Where to Draw the Boundary?}{Eliezer Yudkowsky}{2008}

Someone comes to me and says: ``Long have I pondered the meaning of the word `Art', and at last I've found what seems to me a satisfactory definition: `Art is that which is designed for the purpose of creating a reaction in an audience.''' \nb{The linked comment, by Doug\_S., began ``I've come up with what is, to me, a satisfactory definition''. The words about pondering the meaning of the word were added in the post.}

My answer: a word does not have a meaning floating in the void, which you can discover by finding the right definition. The real challenge is to find which things cluster together, and sometimes which have a common cause. Define ``eluctromugnetism'' to include lightning, compasses and Mesmer's animal magnetism but not light, and you will have trouble asking how it works. That is a wrong word, a cut that fails ``to carve reality along its natural joints.'' Finding where to cut is the problem worthy of a rationalist, and it is not a job for dictionary editors.

Suppose you give me a list of things you call art (a Bach fugue, an Escher print, the Python programming language) and things you do not (a flower, a cross floating in urine, Modern Art), and ask for a simple rule that fits. I reply that the included items inspire similar aesthetic emotions, and that people made them with the intent of producing such emotions. \nb{That is the same kind of answer as the commenter's: a statement of what the things called art have in common. The post gives no reason to treat one as a legitimate guess and the other as a search for an essence.} Is that cheating at Taboo? I would argue that the list of aesthetic emotions is far more compact than the list of artworks.

But my definition is not the point. You could dispute my rule, saying it roughly fits the points but ``is not the true generating distribution.'' Or you could dispute my list and say that Modern Art belongs on it. ``This would mark you as a gullible philistine, but you could argue it.'' Either way the presumption is that some pattern generates the list. \nb{For animal magnetism and for fish the post names what settles the question, a later finding about what the members share. For art it names no such finding.}

Long before anyone knew what electricity and magnetism have in common, you could suspect that animal magnetism did not belong with them. \nb{In 1784 a French royal commission that included Franklin and Lavoisier found no evidence for Mesmer's ``magnetic fluid''.} People once counted dolphins as fish. You could insist that a list cannot be wrong, or you could ``stop playing nitwit games'' and admit that dolphins don't belong. \nb{Linnaeus classed whales as fish in 1735 and moved them to the mammals in 1758.}

So both kinds of definition, the rule and the list, can be wrong. Categorizing is guessing, and you should be able to admit that your definitions can be ``mistaken.''
